\chapter{Imparare musica, diventare risorsa per qualcuno}\label{sec:learning}

Il valore educativo del Portale comincia da una possibilità molto concreta:
una persona che non si considera musicista può produrre un suono, ascoltare
una risposta e scoprire di avere qualcosa da proporre. La città le offre
uno strumento e un interlocutore prima di chiederle un'iscrizione, un titolo
o una prestazione. Questa apertura costituisce una scelta politica del
progetto; la sua efficacia educativa resta da verificare. Un oggetto molto
frequentato potrebbe offrire soprattutto intrattenimento. Un laboratorio
musicalmente riuscito potrebbe raggiungere soltanto chi possiede già
strumenti culturali e tempo libero. La ricerca deve rendere visibili entrambe
le possibilità, senza ridurre il piacere di suonare a un punteggio.

\section{Una teoria del cambiamento esposta alla prova}

Proponiamo una catena di cinque passaggi. L'accesso gratuito e comprensibile
abbassa il costo del primo tentativo; una risposta sonora controllabile rende
leggibile il rapporto tra gesto e risultato; un compito condiviso richiede
ascolto e negoziazione; la restituzione tra partecipanti rende comunicabile
una scoperta; una risorsa riutilizzabile e un appuntamento successivo
permettono di riprenderla. Le risorse necessarie sono uno strumento affidabile,
un ambiente praticabile, tempo di facilitazione e una piccola redazione
dei materiali. Gli esiti attesi sono autonomia musicale circoscritta,
capacità di aiutare un compagno e continuità della pratica. Sono ipotesi
del progetto, non risultati già osservati.

Ogni passaggio può interrompersi. Un'interfaccia appariscente può attirare
senza risultare comprensibile; una coppia può dividersi in esecutore e
spettatore; un tutorial può sostituire l'ascolto; un esercizio salvato può
non essere più aperto. Registrare soltanto accessi o minuti di permanenza
nasconderebbe questi fallimenti. La valutazione deve seguire le transizioni:
quanti iniziano senza aiuto, quanti assumono entrambi i ruoli, che cosa
sanno rifare, che cosa riprendono in seguito. Il diritto di partecipare
rimane indipendente dall'accettazione di essere valutati.

\section{Peeragogy, corpo e attenzione condivisa}

Nel \emph{Peeragogy Handbook}, l'apprendimento e la produzione di conoscenza
tra partecipanti si organizzano attraverso obiettivi e strutture che il
gruppo può modificare \cite{learn_peeragogy}. Nel Portale questa impostazione
diventa una regola operativa: ciascuno deve poter scegliere, provare,
spiegare e correggere una parte del lavoro comune. Essere pari non significa
possedere la stessa competenza. Un pianista esperto può apprendere da chi
propone un ritmo, conosce un repertorio diverso o riconosce una difficoltà
di accesso che il pianista non incontra.

La dimensione corporea va progettata nei gesti: percepire la durata di una
pressione, anticipare un ingresso, dosare l'intensità, lasciare spazio al
suono altrui. Per il prototipo assumiamo che rendere questi rapporti
percepibili favorisca l'esplorazione intenzionale; sarà necessario
osservarlo. Lo schermo può visualizzare l'alternanza dei turni o una
struttura ritmica, ma deve poter essere oscurato. Se il partecipante segue
soltanto un indicatore luminoso e non riconosce più la frase ascoltandola,
il supporto ha prodotto dipendenza dal suggerimento. La prova senza schermo
serve precisamente a distinguere le due situazioni.

Una \emph{scoping review} di 28 studi empirici pubblicati tra il 2006 e il
2023 descrive motivazione, identità e pratiche dell'apprendimento musicale
informale \cite{weatherly2025}. Pubblicata online nel settembre 2024 e nel
fascicolo del 2025, offre una mappa, non una stima unica dell'effetto. Gli
autori discutono limiti di generalizzazione e campioni piccoli; nella
selezione sono esclusi lavori focalizzati sulla musica digitale o
sull'apprendimento virtuale. Non sostiene dunque, da sola, l'efficacia
della componente digitale proposta qui.

Anche il legame tra musica e cooperazione richiede misura. Un esperimento
con bambini di quattro anni rilevò comportamenti di aiuto successivi
all'attività musicale condivisa \cite{learn_kirschner}; un successivo
esperimento sul campo non trovò gli effetti attesi della sincronizzazione
\cite{learn_baier}. Popolazioni, compiti e contesti non coincidono con un
parco urbano intergenerazionale. Assumiamo la cooperazione come pratica
da costruire e osservare, senza promettere che suonare insieme renda
automaticamente più solidale una comunità.

\section{Tre situazioni, tre apprendimenti osservabili}

\paragraph{Una domanda di quattro suoni.}
Due persone dispongono di dodici minuti. La prima inventa una frase di
quattro suoni su cinque tasti disponibili; la seconda la riprende e cambia
un elemento scelto: conclusione, durata o intensità. Dopo tre scambi
invertono i ruoli. Il facilitatore chiede: ``Che cosa hai mantenuto e che
cosa hai cambiato?'' Il materiale iniziale contiene una dimostrazione breve
e un foglio con simboli semplici, senza richiedere lettura del pentagramma.
L'obiettivo è distinguere ripetizione e trasformazione intenzionale.
L'evidenza è una nuova risposta coerente con il vincolo, seguita da una
spiegazione verbale, gestuale o musicale. Se il compagno esegue tutto,
il compito non è completato, anche quando il risultato suona bene.

\paragraph{Tenere il tempo, lasciare entrare l'altro.}
Una coppia lavora per venti minuti su due battute: un ostinato di quattro
pulsazioni e una risposta ritmica. Può scegliere tempo e timbro, usando
una tastiera divisa o due controlli accessibili. Prima prova con un
riferimento regolare, poi lo elimina per quattro ripetizioni. Si ferma,
ascolta e decide una sola correzione. I ruoli cambiano; una registrazione
MIDI locale, soltanto se accettata, permette di confrontare gli attacchi.
L'obiettivo comprende mantenere il riferimento e recuperare dopo una
perdita del tempo. Una ripartenza negoziata vale più della prosecuzione
meccanica di un accompagnamento automatico. La misura temporale accompagna
il giudizio musicale e non penalizza variazioni espressive concordate.

\paragraph{Una cartolina sonora che un'altra persona può suonare.}
In quaranta minuti, tre o quattro partecipanti compongono un frammento
originale di otto battute ispirato a un'esperienza del luogo. La relazione
con il parco può essere una pausa, un'andatura o un racconto: non occorre
registrare passanti o suoni ambientali. Il gruppo sceglie struttura,
contrasti e conclusione, prepara un audio dimostrativo e istruzioni
accessibili. Un secondo gruppo prova il materiale senza spiegazioni
aggiuntive degli autori e segnala dove si perde. La revisione di quelle
istruzioni costituisce parte del lavoro musicale. L'esito è un contributo
trasmissibile, con provenienza e condizioni di riuso controllate; il
consenso alla pubblicazione viene chiesto separatamente dalla partecipazione.

\section{Il lavoro del facilitatore e l'accesso reale}

Il facilitatore prepara materiali con difficoltà graduata, esplicita i
turni, osserva chi non riesce a entrare e interviene prima che l'abilità di
uno diventi il silenzio degli altri. Durante il compito propone una domanda
o una dimostrazione breve; poi restituisce il controllo. Nella scheda di
sessione annota numero e tipo degli aiuti, perché un risultato ottenuto
con cinque interventi non equivale allo stesso risultato autonomo. La
formazione comprende gestione degli errori, ascolto, adattamento dei
controlli e procedure di interruzione. Non basta saper suonare.

Prevediamo una modalità senza testo esteso, istruzioni ascoltabili,
dimostrazioni visive prive di lampeggiamenti, controlli adatti a una sola
mano, volume regolabile e tempi non competitivi. Questi sono requisiti da
provare con persone interessate, non una certificazione di accessibilità.
Ogni adattamento va descritto insieme al compito: una misura basata sulla
velocità delle dita non rappresenta la comprensione di chi usa un
controllo alternativo. Per le attività valutate si concorda una modalità
equivalente rispetto all'obiettivo; quando l'equivalenza non è sostenibile,
si descrive il caso senza confrontare punteggi impropriamente.

L'osservazione pubblica può scoraggiare più della difficoltà tecnica.
Servono prove a volume contenuto, possibilità di ascoltare prima di
esporsi e sessioni senza pubblico raccolto attorno allo strumento.
L'invito deve raggiungere anche chi non frequenta già scuole di musica.
Biblioteche e associazioni potranno aiutare a reclutare pubblici diversi,
se accetteranno di collaborare: non sono partner acquisiti. Nel primo
pilota proponiamo adulti; attività con minori richiedono un protocollo
dedicato e personale preparato, anziché un'estensione automatica dei
risultati degli adulti.

\section{Una memoria comune con diritti leggibili}

La raccomandazione UNESCO sulle OER richiede materiali in pubblico dominio
o con licenze che permettano accesso, riuso, adattamento e redistribuzione
\cite{unesco_oer2019}. Proponiamo per le nuove schede originali CC BY 4.0,
che consente anche adattamenti con attribuzione e indicazione delle
modifiche \cite{learn_ccby}. La concessione è irrevocabile per chi rispetta le condizioni: ritirare
un file dal catalogo non cancella i riusi già autorizzati.
È una proposta editoriale: la licenza non
viene applicata automaticamente a contributi di terzi o alle persone
registrate. Un brano reperibile online non diventa riutilizzabile perché
serve a insegnare musica.

Ogni risorsa avrà identificatore, autore o pseudonimo concordato, data,
obiettivo, prerequisiti, durata, versione, formato modificabile, licenza
e fonti degli eventuali materiali incorporati. Il curatore verifica diritti
e comprensibilità prima della pubblicazione. Se manca una base chiara per
il riuso, il contributo resta fuori dall'archivio pubblico. Si distingue
la composizione dall'esecuzione registrata e dalle istruzioni. Una
raccolta iniziale di sei schede revisionate è preferibile a un flusso
indiscriminato: permette di misurare manutenzione, riusi e correzioni.
L'accesso offline conserva utilità anche senza account o telefono.

\section{Pilota di fattibilità e apprendimento trasferibile}

Dentro la finestra di dodici settimane dello scenario economico A,
proponiamo una raccolta educativa di sei settimane, con 36 adulti organizzati in 18 coppie,
reclutati dichiarando esperienza musicale e disponibilità al richiamo.
Sono dimensioni di fattibilità dettate da tempo e presidio, non derivate
da un calcolo di potenza. Prima del reclutamento si pubblicano protocollo,
rubrica, criteri di esclusione e analisi prevista. Le prove di comprensibilità precedono l’assegnazione e non entrano nel
confronto fra condizioni; il loro numero e i cambiamenti conseguenti
vengono registrati separatamente. Le sessioni valutate mantengono
stabili esercizi e versione del software, annotando ogni modifica urgente.
Una sessione dura circa trenta minuti; il richiamo, dopo 7--14 giorni,
richiede dieci minuti.

Le coppie vengono assegnate casualmente a due condizioni, bilanciando per
quanto possibile l'esperienza iniziale: esplorazione con istruzioni di base
oppure compito strutturato con alternanza e restituzione tra pari.
Entrambe hanno uguali durata, strumento e disponibilità di assistenza
tecnica. L'unità assegnata è la coppia; trattare le 36 persone come
osservazioni indipendenti gonfierebbe l'informazione. Il confronto stima
esplorativamente il valore del pacchetto di facilitazione, non l'effetto
puro di una singola componente né il vantaggio del Portale rispetto a
una tastiera più semplice. Quest’ultimo confronto, necessario prima di
giustificare l’integrazione definitiva, richiede una fase distinta con
stessi compiti e stessa dose di facilitazione nelle due configurazioni. Lo stesso facilitatore opera in entrambe
le condizioni con un copione e tempi registrati. Soltanto dopo il richiamo differito, tutti
possono provare la modalità strutturata, evitando di alterare il confronto
prima della misura di mantenimento. La pratica svolta altrove fra le
visite viene registrata come possibile fonte di contaminazione.

Prima e dopo si richiede una breve frase ritmica; al richiamo si propone
un materiale equivalente ma nuovo. Le prove avvengono senza AI,
accompagnamento automatico, suggerimenti luminosi o aiuto del compagno.
L'AI resta esclusa da questa prima fase per rendere interpretabile il
contributo della struttura sociale. In un'eventuale fase successiva si
potrà confrontare un supporto AI specifico mantenendo queste prove autonome.
Il trasferimento riguarda un compito musicale vicino a quello esercitato;
non autorizza deduzioni su intelligenza, rendimento scolastico o benessere.

\paragraph{Misure e qualità del giudizio.}
Per il ritmo si riporta l'errore assoluto mediano degli attacchi rispetto
al riferimento, normalizzato per la durata della pulsazione, usando un
allineamento dichiarato e annotando note mancanti o aggiunte. La precisione
da sola non misura musicalità. Una rubrica a tre dimensioni valuta
riconoscibilità della frase, trasformazione intenzionale e recupero dopo
interruzione, ciascuna da 0 a 2: assente, parziale, autonomo. Due valutatori
indipendenti, senza conoscere condizione e momento della prova quando
possibile, giudicano i file codificati; si pubblicano accordo esatto,
divergenze e decisioni, prima di qualunque media aggregata.

La componente tra pari si osserva attraverso alternanza effettiva,
richieste di chiarimento e correzioni accolte: descrizioni con esempi,
senza trasformarle in una scala psicologica non validata. Si riportano
anche abbandoni, aiuti, accessi impediti, richiami mancanti e motivi
volontariamente forniti. I risultati comprendono distribuzioni e
incertezza, non soltanto una percentuale migliorata. Novità dell'oggetto,
autoselezione, pratica tra le visite e differenze tra coppie limitano
l'inferenza; un confronto piccolo non dimostra efficacia cittadina.

\paragraph{Regole per proseguire, correggere o fermarsi.}
Le seguenti soglie sono proposte gestionali da concordare prima del
pilota. Interrompere immediatamente una sessione in caso di disagio
espresso, problema di sicurezza o acquisizione di dati non prevista.
Sospendere il protocollo interessato fino alla correzione se un incidente
si ripete o emerge un ostacolo che esclude sistematicamente un pubblico.
Rivedere reclutamento e richiamo se torna meno del 60\% degli iscritti;
non dichiarare allora la ritenzione dell'apprendimento dell'intero gruppo.
Rivedere la rubrica se l'accordo esatto dei valutatori è inferiore al
75\% delle valutazioni dimensionali. Sono segnali di revisione, non soglie
scientifiche universali. Un'estensione richiede fattibilità documentata,
accesso migliorabile e risultati interpretabili; l'assenza di un vantaggio
misurato va pubblicata insieme alle esperienze riuscite.
